\documentclass[11pt]{article}
\usepackage[T1]{fontenc}
\usepackage{lmodern}
\usepackage[utf8]{inputenc}
\usepackage[margin=1in]{geometry}
\usepackage{amsmath,amssymb,amsthm,booktabs,longtable}

\newcommand{\Qamp}{\mathcal Q_{\mathrm{amp}}}
\newcommand{\PhiSrc}{\Phi_{\mathrm{src}}}
\newtheorem{definition}{Definition}
\newtheorem{theorem}{Theorem}
\newtheorem{proposition}{Proposition}
\newtheorem{obstruction}{Scoped obstruction}
\newtheorem{boundary}{Boundary}

\title{P5-T05 Backgrounds, Phases, and Perturbative Modes of the
Cubic-Amplitude Source Dynamics}
\author{Candidate Constructor parent--child synthesis}
\date{2026-07-25}

\begin{document}
\maketitle

\begin{center}
\fbox{\parbox{0.92\linewidth}{
\textbf{Status and authority.} This object is \texttt{draft/control},
\texttt{proposal-only}, and \texttt{source-extension data}.  It analyzes the
exact unchanged P5-T03 cubic-amplitude candidate using the P5-T04 structural
results.  It is not a canonical-ontology candidate, an adopted or rejected
source law, a physical phase, a physical mode count, or a source-to-target
reconstruction.  Current adoption remains blocked while same-milestone
continuation remains open.
}}
\end{center}

\section{Source basis and candidate lock}

Let
\[
 \mathcal A=\mathbb T^4,\qquad
 e(\theta)=\sin\theta_1,\qquad
 q_a=a e,\qquad
 \Qamp=\{q_a:a\in\mathbb R\}.
\]
The exact unchanged candidate is
\begin{equation}
 X_\gamma(q_a)=-\gamma a^3e,\qquad
 \frac{da}{d\lambda}=-\gamma a^3,\qquad
 f^\gamma_\lambda(a)
 =\frac{a}{\sqrt{1+2\gamma a^2\lambda}},
 \label{eq:candidate}
\end{equation}
for \(\gamma>0\) and \(\lambda\geq0\).  P5-T04 proved that this is a
one-global-mode family: independent patch amplitudes cannot be glued inside
\(\Qamp\).  It also established identity-only declared redundancy, the
non-gauge reflection \(a\mapsto-a\), and the admitted tangent multiplier.
This packet neither alters those results nor supplies off-family dynamics.

No target atlas, target metric, exact-GR solution, target covariance,
benchmark fit, generated derivative, registry metadata, role authority, or
validator result is used as a scientific premise.  The planned path
\texttt{research\_control/support\_formalization/} is absent and supplies no
authority.

\section{Background and phase catalog}

\begin{definition}[Source background]
A stationary background for this candidate is a state \(q_{\bar a}\in\Qamp\)
for which \(X_\gamma(q_{\bar a})=0\).  An exact trajectory background is a
solution \(q_{a_b(\lambda)}\) of equation~\eqref{eq:candidate}.  A
\emph{reconstruction-suitable phase} would additionally need source-defined
local perturbative data capable of supporting a controlled coarse-graining
or reconstruction map.  That last property is tested here; it is not assumed.
\end{definition}

\begin{theorem}[Unique stationary background and invariant sectors]
\label{thm:backgrounds}
For \(\gamma>0\), the unique stationary background is
\[
 B_0=q_0.
\]
The three strata
\[
 \Qamp^+=\{q_a:a>0\},\qquad
 \Qamp^0=\{q_0\},\qquad
 \Qamp^-=\{q_a:a<0\}
\]
are forward invariant.  Reflection exchanges \(\Qamp^+\) and
\(\Qamp^-\), fixes \(\Qamp^0\), and the dynamics selects neither nonzero
sign.  Every nonzero trajectory remains in its sign stratum and converges to
\(B_0\).
\end{theorem}

\begin{proof}
The fixed-point equation is \(-\gamma\bar a^3=0\), hence
\(\bar a=0\).  Equation~\eqref{eq:candidate} preserves the sign of \(a\), and
its denominator diverges as \(\lambda\to\infty\) for nonzero \(a\), so
\(a(\lambda)\to0\).  Oddness of the flow gives the reflection statement.
\end{proof}

The nonzero strata are transient invariant sectors, not stationary
thermodynamic or physical phases.  If Borel probabilities are adjoined as
extra structure, the P5-T04 theorem supplies the sole invariant probability
\(\delta_0\); the candidate itself selects no ensemble.

\section{Linearized equations}

Let \(a_b(\lambda)\) be an exact trajectory background with initial amplitude
\(a_{b0}\).  Write
\[
 a(\lambda)=a_b(\lambda)+\epsilon\eta(\lambda)+O(\epsilon^2).
\]
Linearization gives
\begin{equation}
 \frac{d\eta}{d\lambda}
 =-3\gamma a_b(\lambda)^2\eta,
 \qquad
 \eta(\lambda)
 =\frac{\eta_0}
 {(1+2\gamma a_{b0}^2\lambda)^{3/2}}.
 \label{eq:linearized}
\end{equation}
This is the complete linearized equation on the admitted tangent line
\(T_{q_a}\Qamp=\operatorname{span}\{e\}\).

\begin{theorem}[Mode and stability classification]
\label{thm:modes}
For the admitted candidate domain \(\gamma>0\):
\begin{enumerate}
 \item about any nonzero trajectory background, the one admitted tangent mode
 is strictly decaying for positive parameter;
 \item about \(B_0\), the linearized generator is zero, so the one admitted
 tangent mode is neutral at linear order;
 \item \(B_0\) is nevertheless globally Lyapunov stable and globally
 asymptotically attractive under the nonlinear flow;
 \item the attraction is polynomial, not exponential; and
 \item there are no admitted growing or oscillatory modes.
\end{enumerate}
\end{theorem}

\begin{proof}
Equation~\eqref{eq:linearized} is positive and strictly smaller in magnitude
than its initial value when \(a_{b0}\neq0\) and \(\lambda>0\).  At
\(a_{b0}=0\), it reduces to \(\eta' =0\).  For the nonlinear flow,
\[
 |f^\gamma_\lambda(a_0)|\leq |a_0|,\qquad
 \lim_{\lambda\to\infty}f^\gamma_\lambda(a_0)=0,
\]
which proves stability and global attraction.  For \(a_0\neq0\),
\[
 |f^\gamma_\lambda(a_0)|
 \sim (2\gamma\lambda)^{-1/2},
\]
so no uniform positive exponential decay rate exists.
\end{proof}

\begin{boundary}[What is and is not counted]
The candidate has one admitted source-tangent degree of freedom.  Declared
redundancy is identity-only, so there are zero gauge modes.  Hamiltonian
constrained-mode counts are not applicable because no Hamiltonian formalism
exists.  Off-family perturbations are excluded by the current state-space
definition but have no supplied dynamics, so they are not counted as zero
physical or constrained modes.  The one admitted mode is not thereby a
physical mode.
\end{boundary}

\section{Symmetry, order, and parameter controls}

The reflection symmetry is unbroken by the unique stationary background.
The sign strata retain initial-condition labels, but neither is dynamically
preferred and both decay to the symmetric fixed point.  Thus the current
candidate supplies no persistent ordered or symmetry-broken stationary
phase.

The forward parameter and nonsurjective action define a mathematical
source-order orientation only.  They do not establish physical time,
causality, entropy production, or a time-oriented physical phase.

\begin{longtable}{p{0.20\linewidth}p{0.53\linewidth}p{0.19\linewidth}}
\toprule
Parameter branch & Exact result & Disposition\\
\midrule
\endhead
\(\gamma>0\) &
One fixed point, global forward totality, nonexpanding admitted tangent
dynamics, polynomial attraction to zero. &
Admitted candidate branch.\\
\(\gamma=0\) &
Every amplitude is stationary and every tangent mode is neutral. &
Degenerate identity control; excluded by P5-T03 nontriviality.\\
\(\gamma<0\) &
For nonzero \(a_0\), the solution blows up at
\(\lambda_\ast=(2|\gamma|a_0^2)^{-1}\); \(B_0\) is nonlinearly unstable. &
Fails global forward totality.\\
Target-selected background &
A background chosen because it resembles a GR solution or desired metric
imports target information. &
Fail closed: target import.\\
Off-family mode claim &
No update equation exists for perturbations outside
\(\operatorname{span}\{e\}\). &
Fail closed: undefined dynamics.\\
\bottomrule
\end{longtable}

\section{Reconstruction-suitability test}

The unique source-selected stationary background is \(B_0\), but its
linearized operator on the complete admitted tangent space is the zero
operator.  The nonzero exact trajectory backgrounds supply only one global
decaying amplitude variation.  By the P5-T04 patching theorem, no independent
local perturbations occur inside \(\Qamp\).  Consequently the current family
contains no momentum-indexed spectrum, local characteristic operator,
propagating sector, or source-defined resolution hierarchy from which a
continuum reconstruction phase could be selected.

\begin{obstruction}[No nontrivial reconstruction-suitable phase in
\(\Qamp\)]
\label{obs:reconstruction-phase}
\texttt{OBST-P5T05-QAMP-NO-RECONSTRUCTION-PHASE-001}: within the exact
unchanged P5-T03 candidate, no stationary nonzero phase exists, the sole
stationary background has a degenerate zero linearized generator, and the
only admitted nonzero-background perturbation is a single global decaying
amplitude mode.  Together with the preserved patch-amplitude obstruction,
these facts block identification of a nontrivial local
reconstruction-suitable phase in this candidate family.
\end{obstruction}

This is a precise candidate-specific obstruction.  It does not prove that a
conservative multi-mode, locally patchable, coupled, stochastic, or other
source extension cannot supply backgrounds and local modes.  It does not
reject the wider project or establish a no-go theorem.

\section{Distance-to-GR disposition and next obligation}

P5-T05 is complete for the exact candidate: its stationary backgrounds,
trajectory sectors, linearized equations, admitted mode counts, nonlinear
stability, symmetry behavior, and parameter controls are explicit.  The
result narrows the source-ontology burden by recording the exact background
and mode structure, while adding the scoped reconstruction-phase obstruction.

P5-T06 may use this negative datum to attempt a fully typed source-only
coarse-graining or reconstruction map with explicit domain, codomain, scale,
information loss, and error.  It must not erase the obstruction or manufacture
local degrees of freedom, a characteristic cone, or a target metric from the
zero background.

Canonical ontology or source-law adoption, a physical phase or mode count,
physical time, \(M_{\mathrm{src}}\), \(g_{\mathrm{eff}}\), matter coupling,
Einstein equations, exact-GR recovery, benchmark promotion, Gate Chair
closure, proof publication, push, completed derivation, global theory
rejection, and future source-extension impossibility remain blocked.

\end{document}
